\begin{claim}{C.78}{Proved}{Rank-one endpoint images with a controlled joint character and base lift}
Keep the actual disk, proper one-dimensional interior
operator group $Q$, original relative loop, principal
coordinate, and set $A$ of C.76--C.77. For each
component $I$ of $\mathbb R\setminus A$, let $B_I$
be its actual base image, and let $V_I\subset\mathbb R^k$,
$0\leq k\leq2$, be the actual vector image of all
finite-endpoint characters. Retain the strips satisfying
\[
                  B_I\text{ is abelian}
                       \quad\hbox{or}\quad\dim V_I\leq1.
                                                               \tag{C.78.1}
\]
Let $Z_1$ consist of these strips, all lines over $A$,
and the original boundary circle. The original loop
restricted to this closed set has a based contraction
through its actual homeomorphisms, fixing the circle
and commuting with each original interior operator.
Both the contraction and its inverse are jointly
continuous. On every line over $a\in A$ its value
is exactly $(s,a)\mapsto(s+r u_\tau(a),a)$.

The construction includes arbitrary accumulation of
strips at interior lines and at the original circle.
It gives an explicit continuous restriction homomorphism
as in (C.76.4), with $Z_1$ as target subset. When
every strip is retained it contracts the original
loop in its original relative centralizer. For a
proper $Z_1$ no lift through the excluded strips is
asserted. Every excluded strip is a bounded interval
$I=(a,b)$ with nonabelian base image and
\[
                    (c_a,c_b)(G)=\mathbb R^2.
                                                               \tag{C.78.2}
\]

Two exact constructions supply the new assertion.
Compact convex averages of actual source maps give
a continuous increasing base deformation. The C.62
triangle in the scalar-character/base comparison plane
realizes every such value by one compact three-word
family of actual original maps. A second selection
uses compact local pair transports for the exact
Polish path-group actions. No continuous section of
an orbit or character map is assumed.
\end{claim}

\begin{proof}
\emph{The first compact source family.}
Use the original closed relative group $L$ of C.76,
which is the closure of the group generated by the
original loop in its actual closed-disk topology.
Every member has the principal form
\[
                    k(s,t)=(s+u_k(t),\beta_k(t)),
                                                               \tag{C.78.3}
\]
preserves each interval and each line over $A$,
fixes the original circle, and commutes with $Q$.
Its maps and inverse maps are controlled by a
complete compatible metric. Put
$P=\{f_\tau\}\cup\{f_\tau^{-1}\}$.
Construct the compact path-connected family $K$
and its nested compact families $K_n$ exactly as
in (C.77.9), starting with $K_0=P^5$ and using
fifteen short increments at each stage. Their
compact parameter products and uniform limits
take place in the actual group $L$.

The newly retained strips have nonabelian $B_I$
and $\dim V_I=1$. Choose one nonzero endpoint
character $c=c_a$ on such a strip. Every other
finite-endpoint character is a real multiple of
$c$: this is exactly the assertion that the
actual vector image $V_I$ is a line. The relation
holds on $L$ by continuity. The scalar image is
$c(G)=c(L)=\mathbb R$.
For $p=(r,\tau)\in[0,1]^2$, put
\[
       F_n(p)=\{k\in K_n:c(k)=r c(f_\tau)\},\qquad
       F_\infty(p)=\{k\in K:c(k)=r c(f_\tau)\}.
                                                               \tag{C.78.4}
\]
These are nonempty compact sets. Nonemptiness
already follows from the original prefix, since
the real path $c(f_\sigma)$ joins $0$ to $c(f_\tau)$.
All labels in a given set have the prescribed
fraction of \emph{every} finite endpoint character.

\emph{A compact space of positive averages, with its exact topology.}
Choose a countable dense sequence $\varphi_j$ in
the unit ball of $C(K,\mathbb R)$ and define
\[
       \chi(k)=(2^{-j}\varphi_j(k))_{j\geq1}\in\ell^2,
       \qquad
       \mathcal P(F)=\overline{\operatorname{conv}}\chi(F)
                              \quad(F\subset K\text{ compact}).
                                                               \tag{C.78.5}
\]
The map $\chi$ is continuous and injective: the
dense continuous functions separate distinct points.
Its inverse on its image is continuous by compactness.
The closed convex hull $\mathcal P(K)$ is compact.
Indeed, a finite $\epsilon$-net in $\chi(K)$ places
its convex hull within $\epsilon$ of the convex
hull of a finite set. That latter hull is a compact
finite-dimensional polytope. Thus the full hull
is totally bounded; its closure in the complete
Hilbert space is compact. The same proof applies
to every $\mathcal P(F)$.

Every $q\in\mathcal P(K)$ defines a positive
average $\Lambda_q(h)$ of a continuous function
$h\in C(K)$. To give the exact definition, approximate
$q$ by finite convex combinations of the vectors
$\chi(k)$ and take the corresponding averages of
$h(k)$. This has a unique limit: approximate
$h/\max(1,\|h\|_\infty)$ uniformly by one of the
$\varphi_j$, whose average is the specified
$j$th coordinate multiplied by $2^j$.
The uniform approximation bounds the discrepancy
for any two approximating convex combinations.
It also proves independence of the approximating
sequence. The resulting functional is linear,
positive, sends the constant $1$ to $1$, and has
$|\Lambda_q(h)|\leq\|h\|_\infty$.
The same approximation proves joint continuity
in $q$ and in $h$ with its uniform norm.
For $q\in\mathcal P(F)$ one moreover has
\[
              \min_F h\leq\Lambda_q(h)\leq\max_F h.
                                                               \tag{C.78.6}
\]
Thus no representation theorem for measures or
additional selection theorem is needed.

There is a continuous choice
\[
 q_p\in\mathcal P(F_\infty(p)),\qquad
 q_{0,\tau}=\chi(1),\quad q_{1,\tau}=\chi(f_\tau),\quad
 q_{r,0}=q_{r,1}=\chi(1).
                                                               \tag{C.78.7}
\]
Here is its full lower-limit input. If $k\in F_n(p)$,
the correction $r'c(f_{\tau'})-c(k)$ tends to zero.
C.77's short-increment neighborhood statement realizes
it by an actual arbitrarily small increment $e$ in
the next allowed family. The same increment corrects
every label of $F_n(p)$, since those labels have
the same character. For a finite convex combination
of their $\chi$-images, replace every $k$ by $ek$.
Uniform continuity of $\chi$ on the compact source
family makes the new combination arbitrarily close.
Approximation by finite convex combinations therefore
proves
\[
       \mathcal P(F_n(p))\subset
          \liminf_{p'\to p}\mathcal P(F_{n+1}(p')).
                                                               \tag{C.78.8}
\]
Replace these convex sets on the four parameter
faces by the prescribed singletons in (C.78.7).
The lower-limit inclusion survives by continuity
of those prescribed values. Apply the elementary
relative approximation argument of (C.63.22), now
in the Hilbert space: extend one selected vector
locally as a constant, use a partition of unity,
and make successive errors less than $2^{-n}$
with successive differences less than $2^{1-n}$.
Every approximation stays in the convex compact
set $\mathcal P(K)$. Completeness gives a continuous
limit. At each parameter all approximating allowed
sets lie in the compact closed
$\mathcal P(F_\infty(p))$, so the limit lies there;
on the parameter faces the error to the prescribed
singleton tends to zero. This proves (C.78.7).

\emph{The base average and its actual three-word realization.}
On the new strip define
\[
                B_p(t)=\Lambda_{q_p}\bigl(k\mapsto\beta_k(t)\bigr).
                                                               \tag{C.78.9}
\]
It is jointly continuous: on each compact real
coordinate interval, $t\mapsto(k\mapsto\beta_k(t))$
is continuous into $C(K)$ in its uniform norm.
For $t_1<t_2$, compactness gives
\[
       \min_{k\in K}\bigl(\beta_k(t_2)-\beta_k(t_1)\bigr)>0.
                                                               \tag{C.78.10}
\]
Positivity of the average makes $B_p$ strictly
increasing. It is proper onto $I$: for each
$v\in I$, the compact set
$\{\beta_k^{-1}(v):k\in K\}$ is bounded away
from both endpoints of $I$. Beyond its upper
bound every $\beta_k(t)$ exceeds $v$, and below
its lower bound every $\beta_k(t)$ is less than
$v$. The same statements hold for their average.
Thus $B_p$ is an increasing homeomorphism of $I$.
Strict increase and joint continuity give a
jointly continuous inverse, by bracketing its
value between two fixed nearby input points.
The four faces in (C.78.7) give the required
identity and original base maps exactly.

Use the actual comparison-plane action
\[
       \begin{split}
       T_I(k)(z,t)&=(z+c(k),\beta_k(t)),\\
       T_I(k)^{-1}(z,t)&=(z-c(k),\beta_k^{-1}(t)).
       \end{split}
                                                               \tag{C.78.11}
\]
The source $\mathbb R\times I$ is a plane; its
proper real fibers for the target $z$ are the
vertical copies of $I$. Every displayed map
preserves that partition and induces an increasing
translation on its real target. By (C.78.6), the
value $B_p(t)$ lies between the smallest and
largest values $\beta_k(t)$ for $k\in F_\infty(p)$.
Both extrema have actual compact labels in $K$,
and those endpoints have the same $z$ coordinate
$r c(f_\tau)$ in (C.78.11).

Connect their labels by a path in $K$ and apply
the C.62 triangle and zero-set argument in this
comparison plane, using $z$ as the exact real
target. Its proof only requires proper real
fibers and increasing target maps; both have
just been checked. It does not require affine
maps along the fibers. The resulting three-letter
word retains its actual labels in $L$ and realizes
the desired comparison point. Thus, with
\[
                     M_0=K K^{-1}K,
                                                               \tag{C.78.12}
\]
for every $p,t$ there is an actual $w\in M_0$
such that $c(w)=r c(f_\tau)$ and
$\beta_w(t)=B_p(t)$. This is a pointwise source
realization in one compact family. No continuous
choice of $w$ has been asserted.

\emph{All comparison spaces and compact local pair transports.}
For the strips retained already in C.77 use its
base prescription. An abelian base strip uses
the joint character $a_I$ of (C.77.6) and target
space $a_I(G)$ with its translation action.
Its source point is $0$ and desired target is
$r a_I(f_\tau)$; $P^5\subset M_0$ realizes it.
A nonabelian strip with zero endpoint image uses
target space $I$, its actual base action, source
point $t$, and C.27's median target $B_p(t)$;
an original prefix in $P$ realizes that target.
The new rank-one strips use (C.78.11), source
$(0,t)$, and target $(r c(f_\tau),B_p(t))$.
Their compact initial realization is (C.78.12).

There are countably many such spaces, one for
each retained interval. Denote them by $X_i$.
Each is a Euclidean vector space or a product
of a line with an interval, or an interval.
The actual based path group $\mathscr P(L)$
acts transitively on each: the vector action
uses the original character image, the interval
action uses C.75, and the comparison-plane
action uses (C.77.3)--(C.77.4).
All their orbit maps are open by the same
verified Polish orbit theorem. Every action
uses the original endpoint character or base
map; any action kernel remains its literal
kernel in $\mathscr P(L)$.

We require the following exact strengthening
of C.64's compact lifting construction. Given
an identity neighborhood $U$ in $L$, there is
a compact path-connected family $R\subset U$
containing $1$ with this property: for every
$y\in X_i$ and every identity neighborhood $V$
in $L$, there is a neighborhood $N$ of $y$
such that
\[
 \text{for }y_1,y_2\in N\text{ there is }e\in R\cap V
                       \text{ with }e\cdot y_1=y_2.
                                                               \tag{C.78.13}
\]
No continuity of the choice $e$ is part of the
assertion.

To prove it, choose an identity neighborhood
$U_0$ with $U_0U_0^{-1}\subset U$.
Exhaust the disjoint countable union of the
$X_i$ by nested compact sets $E_m$ whose
interiors cover it. Use based path neighborhoods
whose values lie in $U_0$ and within $\epsilon_m$
of $1$, where $\epsilon_m\downarrow0$ in a
compatible metric on $L$. At each point of
$E_m$, openness of its actual orbit map supplies
a small closed coordinate ball $D_{m,j}$ within
the image of that path neighborhood. A finite
collection of their interiors covers $E_m$.
C.64's proved compact-covering lemma gives a
compact family of these small paths whose
endpoints map the chosen center exactly onto
$D_{m,j}$. Let $S$ be the union of all their
path values, together with $1$.

This $S$ is compact and path connected: at each
level there are finitely many compact path-value
families, and the families with $m\to\infty$
converge uniformly to $1$. It lies in $U_0$.
Set $R=SS^{-1}$, so $R$ is compact, path connected,
and contained in $U$. Given $y,V$, take $m$
sufficiently large that $y\in E_m$ and that
products of two level-$m$ values, one inverted,
lie in $V$. Choose a ball $D_{m,j}$ whose
interior contains $y$. For two points in that
interior choose endpoint labels $s_1,s_2$ from
the same compact path family mapping its one
center to those two points. Then
$e=s_2s_1^{-1}\in R\cap V$ maps the first to
the second. This proves (C.78.13) simultaneously
for every comparison space. Both maps and
inverses are controlled in the original group.

\emph{The second compact product and the actual fiber lift.}
Starting with $M_0$ in (C.78.12), choose an
identity neighborhood $U_n$ so small that its
left multiplication changes every previously
formed product by less than $2^{-n}$ in a
complete compatible metric on $L$.
Choose $R_n\subset U_n$ by (C.78.13), and put
\[
                   M_n=R_n\cdots R_1M_0.
                                                               \tag{C.78.14}
\]
Uniform Cauchy convergence on the product of
the compact parameter families gives a compact
path-connected limit family $M\subset L$
containing every $M_n$. Inverse convergence is
included by the actual group topology.

On every retained strip define the fiber interval
using precisely the appropriate target equality:
\[
 A_n(p,t)=\operatorname{conv}
 \{u_k(t):k\in M_n,
          \ k\cdot x_i(t)=y_i(p,t)\}.
                                                               \tag{C.78.15}
\]
Here $(x_i(t),y_i(p,t))$ are the actual source
and target points specified above: respectively
the character-vector pair, the base pair, or
the scalar-character/base pair.
These sets are nonempty and compact by the
initial realizations in $M_0$. Their defining
equality ensures the same actual target base
level $B_p(t)$ in each case.

For a fixed allowed $k\in M_n$ at $(p,t)$,
the two comparison points
$k\cdot x_i(t')$ and $y_i(p',t')$ tend to the
same point as $(p',t')\to(p,t)$.
Property (C.78.13), applied to $R_{n+1}$ and
an arbitrarily small $V$, supplies an actual
$e\in R_{n+1}\cap V$ carrying the first point
to the second. Then $ek\in M_{n+1}$ has exactly
the new target equality. Its fiber displacement
approaches that of $k$ as $V$ decreases to the
identity, uniformly near the specified source
point. Using the same convex coefficient for
two corrected endpoint values proves
\[
                   A_n(p,t)\subset
              \liminf_{(p',t')\to(p,t)}A_{n+1}(p',t').
                                                               \tag{C.78.16}
\]
The four prescribed fiber values are again
(C.76.13), represented by $1$ or $f_\tau$ in
$M_0$. The full relative interval selection
argument of C.76 therefore supplies a continuous
$U(p,t)$ in the limiting compact fiber interval,
with those faces exactly. Use (C.76.18) for the
actual homeomorphism and its displayed inverse.

\emph{Operator transport and every finite endpoint.}
Retain the original component map
$q(s,t)=(\lambda s+\kappa,t+d)$ of C.75,
when it exists. Choose all constructions on
representatives of its orbits of proper intervals.
For a new rank-one strip the chosen endpoint
character transports to $\lambda c$. Its
level sets (C.78.4) for the transported strip
are therefore identical subsets of the same
global $K$. Transport the selected positive
average by using the same vector $q_p$ in
$\mathcal P(K)$. Equation (C.76.17), linearity
and preservation of the constant function by
$\Lambda_{q_p}$ give
$B'_p(t+d)=B_p(t)+d$ exactly.
The allowed actual fiber intervals transform
by multiplication by $\lambda$.
Transport their selection by that multiplier;
the additive $\kappa$ cancels in the actual
commutation equation. For the older retained
strips use precisely the character or median
transport verified in C.77. A nonzero component
power never stabilizes a proper interval.
The interval $I=\mathbb R$ has no endpoint
characters and is one of those older cases;
its possible stabilizer is handled by the
actual line-bundle selection of C.77. Negative
$\lambda$ and every original component label
are consequently retained.

For the new base averages one has the uniform
pointwise bounds
\[
 \begin{split}
 \min_{k\in K}\beta_k(t)&\leq B_p(t)\leq
                          \max_{k\in K}\beta_k(t),\\
 \min_{k\in K}\beta_k^{-1}(v)&\leq B_p^{-1}(v)\leq
                          \max_{k\in K}\beta_k^{-1}(v).
 \end{split}                                                   \tag{C.78.17}
\]
The first is positivity. For the second, an
input below all $\beta_k^{-1}(v)$ has every
$\beta_k$-value below $v$, and an input above
all of them has every value above $v$.
The older base prescriptions and inverses
have the same bounds, since they lie between
the original map and the identity, both in $K$.

Now let retained interval points $t_n$ tend
to $a\in A$, allowing the interval and all
parameters to vary. Compactness of $K$ and
its inverse family, all fixing the base point
$a$, makes both bounds (C.78.17) tend to $a$.
Choose a finite endpoint $a_n$ of the source
interval between $t_n$ and $a$; then $a_n\to a$.
Every label in (C.78.15) has
$u_k(a_n)=r_nu_{\tau_n}(a_n)$: for an abelian
strip this is its vector constraint; for the
new strip it follows from its scalar constraint
and the proportionality of both endpoint
characters; for a zero-character strip both
sides vanish. Uniform continuity over the
one compact family $M$ gives
$u_k(t_n)-u_k(a_n)\to0$, uniformly in its
allowed labels. Convexification preserves
the estimate. This proves the specified
fractional endpoint displacement, including
varying intervals, just as in (C.76.21).
The inverse base bounds in (C.78.17) permit
the identical argument at the inverse input,
with the negative fiber displacement. Thus
both families glue continuously to every
line over $A$.

\emph{The original circle and the exact remaining region.}
Each selected actual fiber value is between
two $M$-values at the same real target level.
The original source-plane C.62 triangle and
inverse bracketing give
\[
       K_{r,\tau}(x)\in M M^{-1}M(x),\qquad
       K_{r,\tau}^{-1}(x)\in M^{-1}M M^{-1}(x).
                                                               \tag{C.78.18}
\]
On lines over $A$, the prescribed scalar
fraction is attained by an original prefix
value and its inverse, so these bounds hold
there too. Both displayed word families are
compact families of the original closed-disk
homeomorphisms and fix the circle pointwise.
They therefore send a sequence tending to
an original boundary point uniformly to that
same point. This proves joint continuity
and inverse continuity at the original
circle, with arbitrarily varying intervals
and unbounded principal coordinates.

The resulting piecewise bijections preserve
all the retained pieces, have the displayed
continuous inverses, and are actual
homeomorphisms of the closed set $Z_1$.
Their compact-parameter evaluations give
continuity in its actual homeomorphism
topology. The restriction subgroup and
homomorphism are precisely those of (C.76.4)
with this named closed subset. On the subgroup
preserving the older subset $Z_{\mathrm{ab}}$,
further restriction is a continuous homomorphism;
every constructed map preserves that subset.
This supplies the exact comparison with the
earlier restricted loop, without asserting
identity of independently chosen fillings.

If a strip remains excluded, its base image
is nonabelian and its endpoint vector image
has dimension two. It must therefore have
two finite endpoints, and their joint image
is all of $\mathbb R^2$, proving (C.78.2).
The scalar-character comparison plane used
in (C.78.11) does not retain both independent
characters of such a strip. No vertical
interval containment in the corresponding
three-dimensional comparison space has
been proved here, and no filling through
that excluded region is inferred.
The other remaining cases of C.12 are
unchanged.
\end{proof}
